\section{Source Integration and Connected Refinement}
\label{app:channelimplementation}

\subsection{Poisson source integration}

For completeness, the inhomogeneous three-dimensional Poisson implementation
used in our engineering benchmark samples the Green source term without a
volume grid.  The following identity also fixes the normalization used in our
implementation.

\begin{proposition}[Unbiased source sample on a three-dimensional WoS ball]
\label{prop:greenballsample}
Let $B_R(x)\subset\Omega$, let $\Theta$ be uniform on $\mathbb S^2$, and let
$T$ be independent of $\Theta$ with density $p_T(t)=6t(1-t)$ on $(0,1)$.
For every measurable source $f$ with $\int_{B_R(x)}G_{B_R(x)}(x,y)|f(y)|\,dy<\infty$,
\begin{equation}
 \int_{B_R(x)}G_{B_R(x)}(x,y)f(y)\,dy
 =\frac{R^2}{6}\,\E[f(x+RT\Theta)],
 \label{eq:greenballsample}
\end{equation}
where $G_{B_R(x)}(x,y)=(4\pi)^{-1}
(\|y-x\|^{-1}-R^{-1})$ is the Dirichlet Green kernel evaluated at the ball
center.
\end{proposition}

\begin{proof}
Denote the left-hand side of \cref{eq:greenballsample} by $I_R(x)$ and write
$y=x+r\theta$.  Spherical coordinates, followed by $r=Rt$, give
\begin{align*}
 I_R(x)
 &=\frac1{4\pi}\int_{\mathbb S^2}\int_0^R
   \left(r-\frac{r^2}{R}\right)f(x+r\theta)\,dr\,d\theta\\
 &=\frac{R^2}{4\pi}\int_{\mathbb S^2}\int_0^1
   t(1-t)f(x+Rt\theta)\,dt\,d\theta.
\end{align*}
Uniform $\Theta$ has surface density $(4\pi)^{-1}$ and $T$ has density
$6t(1-t)$; substituting both densities into the expectation on the right of
\cref{eq:greenballsample} gives exactly the last integral.  Hence one draw of
$(T,\Theta)$ yields an unbiased contribution $R^2f(x+RT\Theta)/6$.
\end{proof}

\subsection{Connected-channel parameterization and refinement}

Topology creation need not end with disconnected cavities.  Let
$q=(q_0,\ldots,q_m)\in(\R^3)^{m+1}$ be control points and let a fixed
subdivision matrix $W\in\R^{L\times(m+1)}$ have nonnegative entries and row
sums one.  Dense centers, the removed channel, and the retained domain are
\begin{equation}
 \begin{aligned}
  s_i(q)&=\sum_{j=0}^{m}W_{ij}q_j,\\
  H(q,r)&=\bigcup_{i=1}^{L}\overline B_r(s_i(q)),
  &\Omega(q,r)&=\Omega\setminus H(q,r).
 \end{aligned}
 \label{eq:channelparameterization}
\end{equation}
Every accepted trial must satisfy $\|s_{i+1}-s_i\|<2r$ for all consecutive
centers, checked again after radius recalibration.  Thus the interiors of
successive balls overlap, making both the removed union and its open interior
path connected.  The
topology stage selects two separated terminals and connects them through a
topological-score-weighted candidate graph; Differential WoS then moves the
interior controls while the endpoints remain fixed.
The current housing experiment keeps every ball strictly inside the housing;
its terminals are interior design points.  It therefore demonstrates an
internally connected cavity, not an inlet-to-outlet conduit or a guarantee of
powder evacuation.  In particular, it would remain a void island under the
exterior-reachability definition of Andrade et al.~\cite{andrade2026connectivity}.
The two constraints must not be treated as interchangeable baselines.

Let $\Sigma$ be the union of sphere-intersection seams on $\partial H$.  For
an exposed patch $\Gamma_i$ and offset $\ell>0$, define the admissible part
of that patch by
\begin{equation}
 \Gamma_{i,\ell}^{\mathrm{safe}}=
 \left\{x\in\Gamma_i:
 \operatorname{dist}(x,\Sigma)>\ell,\;
 (x,x-\ell n(x)]\subset\Omega(q,r)\right\}.
 \label{eq:channelsafepatch}
\end{equation}
The Euclidean seam distance is used here; replacing the unit strip width by
any fixed positive multiple changes only constants.

\begin{proposition}[Differential WoS center derivative]
\label{prop:channel-diffwos}
Fix $r>0$ and a regular channel configuration for which the exposed spherical
patches $\Gamma_i\subset\partial H(q,r)$ are $C^2$ away from $\Sigma$.
Assume the piecewise-smooth solution map is shape differentiable under
independent center translations.  Moreover, for every compact probe set
$K\Subset\Omega(q,r)$, assume there are $C_K,\ell_K>0$ such that, for
$0<\ell\le\ell_K$ and $\xi=X_\tau$,
\begin{align}
 \sup_{y\in K}\E_y\!\left[
  \mathbf 1_{\{I(\xi)=i,\,\xi\notin\Gamma_{i,\ell}^{\mathrm{safe}}\}}
  |\partial_nu(\xi)|\right]&\le C_K\ell,
 \label{eq:weightedseamstrip}\\
 \sup_{x\in\Gamma_{i,\ell}^{\mathrm{safe}}}
 \left|\frac{u(x-\ell n(x))}{\ell}+\partial_nu(x)\right|
 &\le C_K\ell.
 \label{eq:safenormaltaylor}
\end{align}
Let $u$ solve
$-\Delta u=f$ in $\Omega(q,r)$ with
zero Dirichlet data.  For $y\in\Omega(q,r)$, let $X_\tau$ be Brownian exit,
$I(X_\tau)$ its exposed-ball index, and $n$ the outward normal of the retained
domain.  If $U_x$ is an unbiased independent WoS estimate of $u(x)$, then
\begin{equation}
 G_{i,\ell}(y)=
 \mathbf 1_{\{I(X_\tau)=i,\,
 X_\tau\in\Gamma_{i,\ell}^{\mathrm{safe}}\}}
 n(X_\tau)\frac{U_{X_\tau-\ell n(X_\tau)}}{\ell}
 \label{eq:diffwoscenterestimator}
\end{equation}
satisfies, uniformly on compact probe sets,
\begin{equation}
 \E G_{i,\ell}(y)=\nabla_{s_i}u(y)+O(\ell)
 \quad\text{as }\ell\downarrow0.
 \label{eq:diffwoscenterbias}
\end{equation}
For any fixed finite set of interior probes, the weighted estimate has $O(\ell)$ bias for the corresponding weighted derivative sum. To assert this rate for compliance $J(q,r)=\int_{\Omega(q,r)}fu\,dx$, additionally assume that differentiation under the domain integral is valid and that, for all sufficiently small $\ell$, the pointwise estimator error is bounded by $\ell C(y)$ with $\int_{\Omega(q,r)}|f(y)|C(y)\,dy<\infty$. A finite probe rule then adds its separate quadrature error.
\end{proposition}

\begin{proposition}[Control chain rule and length tangency]
\label{prop:channel-chainrule}
Suppose $J$ is differentiable with respect to the dense centers at a regular
configuration, and write $G_i=\nabla_{s_i}J$.  Then
\begin{equation}
 \nabla_{q_j}J=\sum_{i=1}^{L}W_{ij}G_i.
 \label{eq:channelchainrule}
\end{equation}
For the centerline length
$\mathcal L(q)=\sum_{j=0}^{m-1}\|q_{j+1}-q_j\|$, let $c$ be its gradient
with the endpoint blocks set to zero.  Given any endpoint-fixed proposal $a$,
define
\begin{equation}
 d=a-c\frac{\langle c,a\rangle}{\|c\|^2}
 \quad(c\ne0),
 \label{eq:lengthprojection}
\end{equation}
and $d=a$ if $c=0$.  Then the endpoints do not move and
$D\mathcal L(q)[d]=0$.  In particular, let $g$ be the objective gradient with
its endpoint blocks set to zero.  Choosing $a=-g$ gives
$DJ(q)[d]=-\|P_{c^\perp}g\|^2\le0$.  Our implementation first
smooths $a$ and therefore relies on fresh line-search validation, rather than
this last inequality, for accepted descent.
\end{proposition}

The radius of every trial is recalibrated on one fixed Sobol set.  If
$Y_1,\ldots,Y_M$ fill a box of volume $V_B$, define
\begin{equation}
 \widehat V_M(q,r)=\frac{V_B}{M}\sum_{k=1}^{M}
 \mathbf 1_{\{Y_k\in\Omega,\;\min_i\|Y_k-s_i(q)\|\le r\}}.
 \label{eq:qmcchannelvolume}
\end{equation}

\begin{proposition}[Discrete equal-volume calibration]
\label{prop:qmc-volume}
Let $k$ be the nearest integer to $MV_\star/V_B$.  If the $k$th and
$(k+1)$st eligible center distances are distinct, choosing $r$ strictly
between them gives exactly $k$ removed samples and
\begin{equation}
 |\widehat V_M(q,r)-V_\star|\le \frac{V_B}{2M}.
 \label{eq:qmcvolumequantum}
\end{equation}
Thus all compared connected channels have the same discrete volume, up to one
predeclared rounding quantum; this is a statement about the fixed QMC measure,
not an exact analytic union-volume formula.
\end{proposition}

\begin{proposition}[Fresh paired shape acceptance]
\label{prop:fresh-shape}
At shape round $t$, condition on the complete history $\mathcal F_{t-1}$,
including the Differential WoS direction and finitely many trial scales.  Let
$\Delta_{t,k}=J(q_{t,k},r_{t,k})-J(q_t,r_t)$, where every $r_{t,k}$ obeys
\cref{eq:qmcchannelvolume}.  Suppose simultaneous upper bounds satisfy
\begin{equation}
 \Prob\!\left(\Delta_{t,k}\le U_{t,k}\ \text{for every }k
 \mid\mathcal F_{t-1}\right)\ge1-\delta_t.
 \label{eq:shapesimultaneouscoverage}
\end{equation}
If a trial is accepted only when its $U_{t,k}<0$, then
\begin{equation}
 \Prob(\text{some accepted shape step is non-improving})
 \le\sum_t\delta_t.
 \label{eq:shapefamilywise}
\end{equation}
The conclusion is independent of bias in the proposed direction: coverage of
the fresh objective difference, not gradient accuracy, certifies acceptance.
\end{proposition}

